\section{Introduction}
\label{sec:introduction}

A scientific report is not a diary entry, a sales pitch, or a code
comment expanded to ten pages. It is a structured argument: you had a
question or a problem, you did something about it in a principled way,
and you obtained results that a professor, teacher or even a stranger --- your reader --- can
understand, judge, and in principle reproduce. Everything in this guide
follows from that one sentence.

In computer science and engineering, this means your report typically
documents one of the following:
\begin{itemize}[nosep]
      \item an \textbf{experiment} (e.g.\ benchmarking two algorithms, evaluating
            a machine-learning model),
      \item a \textbf{system} you designed and built (e.g.\ a compiler, a web
            service, an embedded controller), or
      \item an \textbf{investigation} of an existing system or dataset.
\end{itemize}
In all three cases, the reader --- your instructor, a fellow student, or
your future self in six months --- should be able to answer three
questions after reading your report: \emph{What did you do? Why did you
do it that way? What did you find, and what does it mean?}

\begin{tipbox}
      Write for a reader who is intelligent and technically competent, but who
      was \emph{not} in the room when you did the work. They do not know what
      you meant to say --- only what you actually wrote.
\end{tipbox}

This document walks through the expected structure of a report
(\cref{sec:structure}), how to write clearly and in an appropriate
academic style (\cref{sec:style}), how to present figures, tables,
algorithms and code (\cref{sec:figures,sec:algorithms}), how to cite
sources correctly (\cref{sec:citations}), a few \LaTeX{}-specific tips
that will make your life easier (\cref{sec:latex-tips}), a list of
mistakes we see every year (\cref{sec:mistakes}), and a checklist to run
through before you submit (\cref{sec:checklist}).
